\providecommand{\Res}[2]{\qty[round-mode=figures,round-precision=3]{\fpeval{#1}}{#2}}
\providecommand{\Val}[1]{\num[round-mode=figures,round-precision=4]{\fpeval{#1}}}
\providecommand{\ResSci}[2]{\qty[round-mode=figures,round-precision=3,exponent-mode=scientific]{\fpeval{#1}}{#2}}
\subsubsection*{Question a) Intensité solaire}
\begin{enumerate}[label=\arabic*)]
    \item \textbf{Variables connues et recherchées :} $S_\oplus = \qty{1361}{W/m^2}$ à \qty{1}{UA}, $d = \qty{\v{dp}}{UA}$ ; variable recherchée : $S$
    \item \textbf{Équations pertinentes :} $I = L/(4\pi d^2)$, donc
    \[ \frac{S}{S_\oplus} = \left(\frac{d_\oplus}{d}\right)^2 \]
    \item \textbf{Isolation de la variable recherchée :}
    \[ S = S_\oplus\left(\frac{\qty{1}{UA}}{d}\right)^2 \]
    \item \textbf{Résolution numérique :}
    \[ S = \frac{1361}{(\num{\v{dp}})^2} = \Res{1361/\v{dp}^2}{W/m^2} \]
    \item \textbf{Réponse :} $S = \mathbf{\Res{1361/\v{dp}^2}{W/m^2}}$.
\end{enumerate}

\subsubsection*{Question b) Température d'équilibre}
\[ T_e = \left(\frac{S(1-a)}{4\sigma}\right)^{1/4} = \left(\frac{\Val{1361/\v{dp}^2}(1 - \num{\v{al}})}{4 \cdot \num{5.67e-8}}\right)^{1/4} = \Res{(1361/\v{dp}^2*(1 - \v{al})/(4*5.67e-8))^0.25}{K} \]
La température d'équilibre est de \textbf{\Res{(1361/\v{dp}^2*(1 - \v{al})/(4*5.67e-8))^0.25}{K}}.

\subsubsection*{Question c) Interprétation}
Une surface beaucoup plus chaude que $T_e$ indique une atmosphère qui piège le rayonnement infrarouge : un fort \textbf{effet de serre} (c'est le cas de Vénus, dont l'atmosphère est très riche en $\mathrm{CO_2}$).
